\section{Introduction}
\label{sec:intro}

The cost of a forward pass in a diffusion transformer
\citep{peebles2023dit}, the backbone of current text-to-image
systems \citep{esser2024scaling,chen2024pixartalpha}, grows
superlinearly with resolution: token count grows quadratically with
edge length, and attention cost grows again in token count. This
cost structure makes the title question worth stating precisely: can
the \emph{gradient} of an adapter be computed on a downscaled latent
as a direct substitute for the native one, with the model, the
objective, and the $\sigma$ distribution unchanged and only
the spatial grid changed on some steps?

There is a well-studied reason to expect an affirmative answer
\emph{at high noise}. As the
noise level $\sigma$ grows, per-frequency signal-to-noise falls, and
high spatial frequencies are the first to be masked by the noise, so
that above a spectral crossover a downscaled latent carries
\emph{almost} the full surviving signal. Studies of this coupling, however, are
concentrated on the \emph{inference} side: scale-wise
distillation \citep{starodubcev2025swd}, spectral progressive
diffusion \citep{xiao2026spd}, and spectrally-guided noise schedules
\citep{esteves2026spectral} run high-noise \emph{sampling} steps at
reduced resolution and judge the premise by the generated output.
Where resolution is reduced during training
\citep{jin2024pyramidal,karras2018progressive,chen2024pixartalpha,hoogeboom2023simple},
the low-resolution stage is given its own, modified objective, and
success is again judged from the final samples. By the spectral
premise, substitution is safe \emph{whenever the noise has masked the
frequencies the coarse grid loses}.

This evidence, however, leaves the question open at \emph{training}
time: whether a downscaled step delivers the native gradient
direction under an unchanged objective. Output-level success observes
the substitution only through the confounds of a full training or
sampling pipeline, and a small difference in a quality score does not
imply a small difference in what a training step \emph{learns}. What
a training step consumes is the expected adapter gradient, and
substituting a downscaled, re-encoded latent for the native one at
training time perturbs more than the network's input: the regression
target carries the clean image at unit weight at every noise level,
and the compute graph itself changes with the grid. We therefore
decompose the perturbation at its source, into a data branch and a
graph branch --- and, measured, the two branches are individually
large and strongly anti-aligned, so that the gap between the
downscaled and native gradients is the residual of a
near-cancellation. On that residual we state an effective two-term
law, a noise-dependent term set by the downscale \emph{ratio}
and a $\sigma$-independent term set by the target grid's \emph{token
count}, and measure each directly.

The contributions of this paper are as follows:

\begin{enumerate}
\item \textbf{Two accounts of the gap, and a debiased measurement
that scores them} (\S\ref{sec:theory}, \S\ref{sec:evidence}). We
define the substitution gap at the gradient level and state two
accounts of it in the same units. The first is the gap
implied by the existing spectral argument, taken in its strongest
tolerance-parameterized form (SPD, \citealp{xiao2026spd}); it
depends on the downscale \emph{ratio} alone by construction. The
second, the \emph{gradient-perturbation} account, treats the
substitution as a perturbation of
both factors of $g = J^{\!\top}r$. Measured, the two perturbation
branches nearly cancel --- strongly anti-aligned at every route and
noise level, an anti-alignment that is uniform across depth and
module type and invariant to the adapter operating point on the
measured checkpoint, and that replicates on two further adapters of
the same base model on the $768$ route --- and on the residual of
that cancellation the account reduces to a route amplitude on a
universal mismatch curve \emph{plus} a $\sigma$-independent term set
by the target grid's token count. The two differ structurally before
any measurement (\S\ref{sec:spectral}), and the resulting
$(\text{route}, \sigma)$ map (\S\ref{sec:map}), read with each
route's finite-sample bias estimated and subtracted and against a
pre-specified margin, decides between them: it uncovers on the
$1024{\to}768$ route a window ($0.65 < \sigma < 0.95$) that the
spectral account does not predict at any tolerance. Scored head to
head on the same curves, the spectral account errs at $0.147$--$0.355$
RMSE, whereas ours predicts held-out routes at ${\sim}0.07$--$0.09$,
and the token-count law correctly predicts the $\sigma$-independent
term on both held-out routes.
\item \textbf{Selective low-resolution training} (\S\ref{sec:practical}).
Restricting the substitution to the measured-safe $1024{\to}896$
route above a one-sided gate ($\sigma > 0.5$) gives
an opt-in $\sigma$-conditional trainer that reduces the forward-compute
footprint by $15.1\%$ and training time by $14.6\%$ at a fixed step
budget, with a weight-space endpoint within kernel-noise reach of a
native retrain, rising to endpoint cosine $0.75$ when the
substitution is scheduled late at proportionally reduced saving
(Table~\ref{tab:yardstick}), the low-excess $1024{\to}768$ window
stacking onto the late rule at no endpoint cost; a visual
counterpart (Fig.~\ref{fig:kaaiyuki}) confirms the endpoint renders
stay close to native.
\end{enumerate}

\section{Related work}
\label{sec:related}

\paragraph{Flow matching and DiT.}
Flow matching \citep{lipman2023flow} regresses a velocity field along
prescribed probability paths between noise and data; rectified flow
\citep{liu2023rectified} takes the paths to be straight lines, giving
the linear noising $z_\sigma = (1-\sigma)x + \sigma\epsilon$ and the
constant target $v = \epsilon - x$ used throughout, the training
objective of current large text-to-image systems
\citep{esser2024scaling}. The Diffusion Transformer
\citep{peebles2023dit} replaces the U-Net denoiser with a transformer
over patchified latent tokens, so a spatial grid enters the network
only as a token sequence: token count grows quadratically with edge
length and attention cost quadratically again in token count, the cost
structure that makes resolution the expensive axis. Spatial position is
now widely carried by a grid-dependent coordinate system, most commonly
rotary embeddings \citep{su2024roformer}.

\paragraph{Scale-wise and progressive-resolution diffusion.}
Progressive growing \citep{karras2018progressive} and resolution
curricula \citep{chen2024pixartalpha,hoogeboom2023simple} train at
increasing resolution as a schedule, each low-resolution stage carrying
its own stage-specific objective. A more recent family ties resolution
to the noise level itself, on the spectral premise above, and has so
far been studied on the \emph{inference} side. Scale-wise distillation
(SwD) \citep{starodubcev2025swd} states the claim directly (above the
crossover, the noised downscaled latent is close in distribution to the
downscaled noised latent) and samples early steps at reduced
resolution. Spectral Progressive Diffusion (SPD) \citep{xiao2026spd}
gives the argument its sharpest available form: a
tolerance-parameterized per-frequency activation time under a diagonal
Gaussian spectral model, a criterion on the per-band output of the
\emph{Bayes-optimal predictor} rather than merely on input SNR, from
which an inference-time resolution schedule follows. Spectrally-guided
noise schedules \citep{esteves2026spectral} apply the same reasoning
to schedule design; where the premise does enter training, as in
pyramidal flow matching \citep{jin2024pyramidal}, the objective is
restaged per pyramid level rather than left native.

\paragraph{Positional extension.}
Position interpolation \citep{chen2023pi}, banded/frequency-selective
variants (YaRN, \citealp{peng2024yarn}), and diffusion-specific dynamic
position extrapolation \citep{issachar2025dype,zhao2026sigma} rescale
rotary coordinates to evaluate a model at unseen grid sizes. We use
exact positional interpolation as a \emph{causal intervention} that
matches the downscaled grid's relative phase geometry to native, in
order to decompose the resolution-sensitivity floor, and adopt the
$\sigma$-gated band schedule of \citet{zhao2026sigma} as a
training-time refinement. On the inference side,
\citet{wu2025phasealigned} measure a content-free rotary attention
bias and show position-interpolated rescaling samples its sinusoidal
phase filter at the wrong phases, ranking interpolation worst among
extension schemes at the sample level; their setting is
mixed-resolution attention within one forward pass, whereas our
substitution is a whole-grid change at training time (no mixed-scale
attention arises) and our attribution is at the training-gradient
level.

\paragraph{Noise levels as tasks.}
A gradient-level literature treats the timesteps of diffusion
training as tasks and measures their pairwise interference: task
affinity, the cosine between per-timestep loss gradients, decays
smoothly as the noise-level gap grows \citep{go2023ant} and turns
negative between distant timesteps \citep{ma2025deme}, motivating
interval clustering and decouple-then-merge training. Those studies
compare full gradients \emph{across} noise levels; our object is
different --- the product-rule split of one gradient's \emph{change}
under input downscaling at a fixed noise level --- and the
near-cancellation central here has no counterpart in that setting.
We cite them as the precedent for smooth, $\sigma$-indexed
gradient-direction structure in these models.
